\documentclass[UTF8,a4paper,11pt,openany]{ctexbook}
\usepackage{../common/statlearnbook}
\SLSetBookMetadata{第 5 册}{统计学习方法讲义 v2.6.0}{采样方法、主题模型与图排序}
\begin{document}
\setcounter{page}{805}
\setcounter{chapter}{33}
\setcounter{figure}{4}
\pagestyle{headings}
\chaptermark{Gibbs 抽样}
\sectionmark{与 Metropolis--Hastings 的关系}
因为$y_{-j}=x_{-j}$，所以精确满条件提议必接受。图\ref{fig:V5-C04-gibbs-vs-mh}同时说明：若只近似满条件或改用其他提议，接受率不再自动等于$1$，此时必须恢复\MetropolisHastings{}校正。
% v2.6.0 figure UID: FIG-P638-01
% Semantic lock: four nodes, four draw paths, formulas and directions unchanged.
\tikzset{
  slfig-FIG-P638-01/.style={font=\fontsize{9.2pt}{11.0pt}\selectfont,every path/.append style={line cap=round,line join=round}},
  slfig-P638-eq/.style={minimum width=36mm,minimum height=13mm,inner xsep=1.5mm,inner ysep=1.6mm,outer sep=1.4pt,fill=white},
  slfig-P638-flow/.style={-{Stealth[length=2.0mm]},draw=SLBlue,line width=.95pt},
  slfig-P638-separator/.style={draw=SLRuleGray,line width=.50pt},
  slfig-P638-warn/.style={-{Stealth[length=1.8mm]},draw=SLRed,line width=.72pt,densely dashed},
  slfig-P638-exception/.style={draw=SLRed,rounded corners=2pt,fill=SLRedBG,line width=.72pt,minimum width=78mm,minimum height=13mm,inner xsep=5pt,inner ysep=4pt,outer sep=.8pt}
}
\begin{figure}[htbp]
\centering
\begin{tikzpicture}[slfig-FIG-P638-01,>=Stealth,every node/.style={font=\fontsize{9.2pt}{11.0pt}\selectfont,align=center}]
  \node[slfig-P638-eq] (q) at (-4.2,0) {{\color{SLBlue}\bfseries 1}\quad 精确满条件提议\\$q_j=\pi(x_j\mid x_{-j})$};
  \node[slfig-P638-eq] (r) at (0,0) {{\color{SLBlue}\bfseries 2}\quad MH 比值逐项抵消\\$R=\dfrac{\pi(y)\pi_j(x_j\mid x_{-j})}{\pi(x)\pi_j(y_j\mid x_{-j})}=1$};
  \node[slfig-P638-eq] (a) at (4.2,0) {{\color{SLBlue}\bfseries 3}\quad $\alpha=1$\\直接接受 $x_j\leftarrow y_j$};
  \draw[slfig-P638-flow] (q)--(r)--(a);
  \draw[slfig-P638-separator] (-6.1,-.85)--(6.1,-.85);
  \node[slfig-P638-exception] (w) at (0,-2.0)
    {近似满条件 / 其他提议 $q_j$\\恢复 MH 接受率校正；拒绝时保持 $x_j$（自环）};
  \draw[slfig-P638-warn] (q.south)--(w.north west);
  \draw[slfig-P638-warn] (a.south)--(w.north east);
\end{tikzpicture}
\captionsetup{width=.94\linewidth}
\caption{精确满条件分布作为单坐标提议时，正流分支的Metropolis--Hastings比值恒为1，Gibbs更新无需拒绝；若提议只近似满条件或改用其他分布，则必须恢复接受率校正和拒绝自环}
\label{fig:V5-C04-gibbs-vs-mh}
\end{figure}

\FloatBarrier
\textbf{优化解释及不适用边界。}标准Gibbs抽样没有要逐轮单调下降的目标函数，因此“优化收敛率”“局部最优点”和“每步改进损失”不适用于本算法。
\end{document}
